\subsection{A quartic auxiliary ordinary seed at the two target weights}
\label{sec:quartic-auxiliary-seed-research}

This construction repairs the specific constant-term obstruction for
the quartic character of conductor five.  It produces an actual cusp
form and its exact local Galois types.  Its extra inertia is killed
over a specified cyclic field.  It does not identify this auxiliary
rational representation with an original deformation point, and
does not transfer the earlier quadratic class certificates.

Put \(p=269\), \(q=3229\), \(K=\mathbb Q(\sqrt{-3})\), and
\(\eta=\chi_{-3}\).  Choose the quartic character \(\psi\) of
conductor five with \(\psi(2)=i\), where \(i^2=-1\) and
\(i\equiv82\pmod p\) in \(\mathbb Z_p\).  Set
\[
 \xi_0=\eta\psi^{-2}=\chi_{-15},\qquad
 \theta(6)=\zeta_{269},\qquad
 \xi=\xi_0\theta,\qquad N=15q.
 \tag{QAS.1}
\]
Here \(\theta\) is the primitive character modulo \(q\) from
AC-G119.  Its square root is \(\theta^{1/2}=\theta^{135}\).
Write \(\Psi=\operatorname{DirG}(\psi)\) and
\(\Theta=\operatorname{DirG}(\theta)\), using arithmetic
Frobenius for these finite avatars and for \(c_p\).
The cohomological representation \(R_g\) has the native geometric
Frobenius convention
\[
 \det(1-XR_g(\operatorname{Fr}^{\rm geom}_r))
   =1-a_r(g)X+\xi(r)r^{k-1}X^2,
 \qquad \det R_g=\operatorname{DirG}(\xi)^{-1}c_p^{1-k}.
 \tag{QAS.2}
\]
The first identity is stated for rational primes \(r\nmid15qp\);
at the coefficient prime \(p\), the ordinary argument uses
crystalline Frobenius instead of an unramified Galois Frobenius.

\begin{theorem}[AC-G127-quartic-auxiliary: actual quartic auxiliary cusp forms]
\label{thm:G127-quartic-auxiliary}
For every odd integer \(k>2\), there is a normalized cuspidal
newform of exact conductor
\[
                  g\in S_k(15q,\xi_0\theta)
 \tag{QAS.3}
\]
and a place of its coefficient field over \(269\) for which
\[
 \begin{aligned}
 a_r(g)&\equiv1+\xi_0(r)r^{k-1}\pmod\varpi
                                      &&(r\nmid15q),\\
 a_3(g)&\equiv a_5(g)\equiv a_q(g)\equiv1\pmod\varpi.
 \end{aligned}
 \tag{QAS.4}
\]
In particular it is ordinary at \(269\).

For \(e=133\) or \(267\), take \(k=2e+1=267\) or \(535\).
The primitive automorphic twist
\[
       f_0=g\otimes(\eta\psi^{-1}\theta^{-1/2})
 \tag{QAS.5}
\]
has nebentype \(\eta\), exact conductor \(75q^2\), and is
ordinary at \(269\).  Its normalized representation
\[
 V=R_{f_0}c_p^e=R_g\eta\Psi\Theta^{1/2}c_p^e
 \tag{QAS.6}
\]
satisfies
\[
 \det V=\eta,\qquad
 \overline V^{\rm ss}
   =\eta\overline\Psi\,\overline c_p^e
           \oplus\overline\Psi^{-1}\overline c_p^{-e}.
 \tag{QAS.7}
\]
It is absolutely irreducible over \(\mathbb Q\) and over \(K\).
At \(p\), it is crystalline ordinary with weights \((-e,e)\)
and the corresponding cyclotomic inertia characters.
Away from \(p\), its exact finite inertia is
\[
 \begin{array}{c|c}
 \text{place}&\text{inertia of }V\\ \hline
 3&1\oplus\eta\\
 5&\Psi\oplus\Psi^{-1}\\
 q&\Theta^{1/2}\oplus\Theta^{-1/2}\\
 \text{every other finite place}&1\oplus1,
 \end{array}
 \tag{QAS.8}
\]
and its monodromy is zero everywhere away from \(p\).

The totally real cyclic field \(T/\mathbb Q\) cut out by
\(\theta\), of degree \(269\), kills precisely the extra
\(q\)-inertia in (QAS.8).  The restrictions of \(V\) to
\(T\) and \(KT\) remain absolutely irreducible.
After a finite value extension, \(V\) has a stable lattice with
exactly the split reduction (QAS.7), and can be written over the
finite complete local test algebra
\(A=\mathbb Z_p+\mathfrak m_{E'}\) with complex conjugation
\(\operatorname{diag}(-1,1)\).
The original coefficient ring \(\mathcal O_0=\mathbb Z_p\)
containing \(i\) is not replaced: the finite value extension
and its place are part of the existence construction.
No original \texttt{IsPt} assertion is made before this
extra inertia is killed.
\end{theorem}

\begin{proof}
\proofstep{1. A genuine integral Eisenstein congruence at all cusps.}
The exact AC-G119 integer certificate gives
\(q-1=12p\), primality of \(q\), and primitive root \(6\).
Here \(q\equiv4\pmod{15}\), not \(1\pmod{15}\), but
\(\xi_0(q)=1\); also \(q\equiv1\pmod p\).
The character \(\xi\) is primitive odd of conductor \(N=15q\).
Choose a finite value extension containing its values and the
prime-to-\(p\) roots of unity for full level \(4N\).
Write \(\mathcal O,\varpi,k_0\) for this extension's valuation
ring, uniformizer and residue field.  Consider the normalized
Eisenstein series
\[
 E_k(1,\xi)=\frac{L(1-k,\xi)}2+
       \sum_{n\ge1}\left(\sum_{d\mid n}\xi(d)d^{k-1}\right)Q^n.
 \tag{QAS.9}
\]
The Fourier variable is \(Q\), distinct from the rational prime.

At the fixed tame level \(N\), Katz's Eisenstein measure gives
the integral moment \(2(1-2^k\xi(2))E_k(1,\xi)\).
This is Theorem 3.3.3 with its exact moment normalization
(3.3.1)--(3.3.2).  Replace \(\xi\) by the function
\(\xi_{00}(a)=\xi_0(a)\) when \(q\nmid a\), and zero otherwise.
Since \(\theta\equiv1\pmod\varpi\), these functions are
congruent on the entire domain \(\mathbb Z/N\mathbb Z\).
The two smoothing factors reduce to \(1-2^k\), because
\(\xi_0(2)=1\).  This is a unit for every odd \(k\):
\(2^{134}=-1\) modulo \(269\) makes \(2^k=1\) impossible
for odd \(k\).  Consequently
\[
 \frac{L(1-k,\xi)}2\equiv
 \frac{L(1-k,\xi_{00})}2
  =(1-\xi_0(q)q^{k-1})\frac{L(1-k,\xi_0)}2
  \equiv0\pmod\varpi.
 \tag{QAS.10}
\]
The same smoothing argument at level fifteen proves integrality
of the last value.  No division by \(k\) modulo \(269\) is used,
so this argument also covers weights divisible by \(269\).
The unit primitive constants at weights \(267,535\) in the
quartic-fork diagnostic are therefore compatible with (QAS.10).

The q-expansion principle makes (QAS.9) integral on the fine
curve below.  The full Dasgupta--Kakde Theorem 4.5(1), applied
with primitive second character \(\xi\) and first character one,
gives zero constant except over \(C_\infty(N,N)\), and there a
root-of-unity multiple of \(L(1-k,\xi)/2\).
The equal Gauss sums cancel, the conductor ratio is one, and
the additional Euler-product set is empty.
All normalizations are prime-to-\(p\) factors.
Thus the reduction \(\overline E\) vanishes at every cusp and
is nonzero because its first Fourier coefficient is one.

\proofstep{2. Integral descent in the order-269 character line.}
Use the compact fine curve \(Y\) for
\(\Gamma_1(N)\cap\Gamma(4)\), on a fixed full-level-four
pairing component.  As in the complete AC-G119 construction,
start from full level \(4N\) and quotient by the upper
unipotent groups at \(3,5,q\).  Their order is \(N\), a unit
at \(p\).  The tame Tate-parameter quotient at a cusp is
\(z\mapsto z^d\) with \(d\mid N\), so the quotient curve is
smooth proper with finite etale reduced cusp divisor \(C_Y\),
and
\[
                   \omega_Y^2=\Omega^1_{Y/\mathcal O}(C_Y).
 \tag{QAS.11}
\]
Let \(D=(\mathbb Z/N\mathbb Z)^\times\) act by diamonds.
Its Sylow \(p\)-subgroup \(P\) has order \(p\), since
\(|D|=8(q-1)=96p\).
It acts freely on the open curve: a fixed point would give
an elliptic-curve automorphism fixing full level four and
multiplying its order-\(N\) point nontrivially, which is impossible.

At a cusp represented by a primitive column \((a,c)^t\), the
parabolic stabilizer, up to sign, is
\[
 \begin{pmatrix}1-acn&a^2n\\-c^2n&1+acn\end{pmatrix}.
 \tag{QAS.12}
\]
Membership in \(\Gamma_0(N)\) forces \(N\mid c^2n\).
Squarefreeness of \(N=3\cdot5\cdot q\) implies \(N\mid cn\),
so both diamond entries are one modulo \(N\), up to the
harmless sign.  No nontrivial element of \(P\) fixes a cusp.
The prime-to-\(p\) full-level cusp labels are finite etale;
the same labels and stabilizers apply to every geometric fiber.
Thus \(Y\to X=Y/P\) is an actual finite etale torsor on the
entire compact curve, by Stacks tag 07S7, including the special
fiber and boundary.  Properness, smoothness and (QAS.11) descend.

Descend the trivial line with \(P\)-character \(\theta^{-1}\)
to a line \(\mathcal L_\theta\) on \(X\), with the convention
that its sections pull back to the \(\theta\) diamond space.
Then \(\mathcal L_\theta^p\simeq\mathcal O_X\), its geometric
fiber degree is zero, and its reduction is canonically trivial.
For \(k>2\), Serre duality gives
\[
 H^1(X_{k_0},\omega_X^k\mathcal L_\theta(-C_X))^\vee
      =H^0(X_{k_0},\omega_X^{2-k}\mathcal L_\theta^{-1})=0,
 \tag{QAS.13}
\]
since \(\deg\omega_X>0\) on every component.  Coherent
cohomology, reduction and Nakayama therefore give a surjection
from integral cusp sections in this character line to its
special-fiber cusp sections.  Lift the descent of \(\overline E\).

Write \(D=P\times D'\), where \(|D'|=96\).
The character \(\theta\) is trivial on \(D'\).
After pullback, project to the \(\xi_0\) diamond space of
\(D'\), then average the commuting full-level-four group
\(\mathrm{SL}_2(\mathbb Z/4\mathbb Z)\), of order \(48\).
Both projectors fix \(\overline E\), and their denominators
are units.  This produces its lift in the exact integral
\(\xi_0\theta\)-character cusp space at level \(N\).
No averaging over \(P\) or division by \(269\) occurs.

\proofstep{3. Eigenvalues, exact conductor and local types.}
The cusp module is finite free over \(\mathcal O\).
Its commuting integral Hecke operators include all good
\(T_r\), in particular \(T_p\), and \(U_3,U_5,U_q\).
The divisor sum in (QAS.9) makes \(\overline E\) a common
eigenvector with exactly (QAS.4), including \(T_p=1\) modulo
\(\varpi\).  The complete Deligne--Serre Lemma 6.11 lifts
these eigenvalues after a finite coefficient extension.
The generated Hecke algebra is finite over \(\mathcal O\),
so the infinite list causes no finiteness problem.
The first coefficient can be normalized to one by
\(a_1(T_n f)=a_n(f)\).  Classical newform decomposition
then supplies a primitive cusp form whose conductor divides
\(N\).  Its primitive nebentype has conductor \(N\), so the
conductor equals \(N\); there is no lower-level oldform ambiguity.
The eigenvalues are algebraic, and their selected \(p\)-adic
embedding supplies the stated coefficient place.

Retain the usual classical attached-representation, purity,
local--global compatibility and ordinary crystalline inputs,
with convention (QAS.2), as in AC-G118--AC-G119.
At each of \(3,5,q\), conductor exponent one and determinant
conductor exponent one force a Weil--Deligne sum of one
unramified and one conductor-one character, with \(N=0\).
An irreducible rank-two Weil representation has conductor
at least two; nonzero monodromy makes the inertia characters
equal and cannot give this conductor/determinant pair.
Thus the respective inertias of \(R_g\) are
\(1\oplus\eta\), \(1\oplus\Psi^2\), and
\(1\oplus\Theta^{-1}\).

At \(p\), the prime-to-\(p\) level and unit \(a_p\) give
crystalline slopes \(0,2e\) and the ordinary filtration.
The attached geometric-Frobenius congruence gives
\(\overline R_g^{\rm ss}=1\oplus
 \operatorname{DirG}(\xi_0)^{-1}\bar c_p^{-2e}\).
The inverse finite twist in (QAS.5) produces exactly (QAS.6).
Since \(\Psi^4=1\), its determinant and reduced characters
are (QAS.7).  At five the twist changes
\(1,\Psi^2\) to \(\Psi,\Psi^3=\Psi^{-1}\), both of
conductor exponent one.  At \(q\) it gives the two nontrivial
square-root characters in (QAS.8); at three it interchanges
\(1,\eta\).  The new conductor is therefore \(3\cdot5^2q^2\).
All other places are unramified, including two, and the
\(p\)-adic twist before \(c_p^e\) is unramified.

Absolute irreducibility over \(\mathbb Q\) is the retained
cuspidal attached-representation theorem.
Reducibility over \(K\) would give the quadratic self-twist
\(V\simeq V\eta\).  At the inert prime \(p\), \(\eta\)
is the unramified character with value \(-1\).
The two distinct crystalline slopes force an isomorphism to
preserve the slope-zero line before the common twist.
Its nonzero Frobenius eigenvalue cannot equal its negative.
Thus this self-twist is impossible, proving irreducibility on
\(G_K\).  For the cyclic degree-\(p\) restrictions, Clifford
semisimplicity and Frobenius reciprocity give the endomorphism
sum over self-twists.  A nontrivial such self-twist would have
\(\nu^2=1\) by determinants, impossible for order \(269\).
This proves the assertions over \(T\) and \(KT\).

For the lattice assertion, begin with a stable lattice over a
finite value DVR.  Its reduction has a composition series with
the two distinct characters in (QAS.7), both valued in
\(\mathbb F_p\).  In a suitable basis its matrix entries satisfy
\(b\in\mathcal O_E\) and \(c\in\varpi_E\mathcal O_E\).
After a ramified quadratic value extension with
\(\varpi_E\) a unit times \(\pi'^2\), conjugation by
\(\operatorname{diag}(\pi',1)\) replaces the off-diagonal
entries by \(\pi'b,c/\pi'\), both divisible by \(\pi'\).
Thus the reduction is split.  Since complex conjugation has
square one and reduction \(\operatorname{diag}(-1,1)\), its
two idempotents \((1\mp V(c))/2\) give a basis congruent to
this one with exactly that diagonal action.  All matrix entries
and inverse entries belong to
\(A=\mathbb Z_p+\mathfrak m_{E'}\).  This is a finite complete
local \(\mathbb Z_p\)-algebra with residue \(\mathbb F_p\);
its topology is the induced valuation topology.  The resulting
representation is continuous over the unchanged original
coefficient ring.  This does not assert the original away-kernel
condition at the auxiliary prime.

Finally the exact AC-G119 computation
\(\theta(p)=\zeta_{269}^{175}\) is unchanged.  The field
\(T\) is totally real, ramified only at \(q\), and unramified
of local degree \(269\) at \(p\).
Its restriction kills (QAS.8)'s \(q\)-inertia.
The residual character is unramified at \(q\), so before
that restriction the nontrivial inertia violates the original
residual-kernel \texttt{Kill} condition.
For both chosen \(e\), \(\psi(-1)=(-1)^e\) and
\(268\nmid2e\); these checks do not remove that failure.
\end{proof}

\paragraph{Exact changed Kummer sectors; no old class certificate transfer.}
For the adjacent residual coefficient
\(\bar\alpha_j=\bar\chi^{-j}\bar\Psi^{-1}\bar c_p^{-e}\), its Kummer
untwist is \(\bar\alpha_j/\bar c_p\).
It has splitting field
\[
 M_d=K(\sqrt[3]{269},F_d),\qquad
 F_d=\mathbb Q(w_d),\quad w_d^4+5d w_d^2+5d^2=0,
 \qquad d=1\text{ or }269,
 \tag{QAS.14}
\]
for \(e=267\) and \(133\), respectively.
Indeed \(F_1=\mathbb Q(\zeta_5)\), with
\(w_1=\zeta_5-\zeta_5^{-1}\) and
\(w_1^2=-(5+\sqrt5)/2\).  When
\(\mathbb Q(\sqrt d)\) is linearly disjoint from \(F_1\), as
for \(d=269\), the element \(w_d=\sqrt d\,w_1\) in their
compositum is fixed by the antidiagonal pair of quadratic
involutions.  It generates the cyclic quartic field cut out by
\(\psi\chi_d\), or its inverse, and contains \(\sqrt5\).
If \(d\in\mathbb Q^{\times2}\), the field is simply \(F_1\).
If \(d\in5\mathbb Q^{\times2}\), then
\(\psi\chi_5=\psi^{-1}\), so the same field \(F_1\) results;
there is no independent quadratic involution in that case.
The polynomial in (QAS.14) remains valid in both exceptional cases.
The quadratic character \(\bar c_p^{134}\) in the second branch
is \(\chi_{269}\), as proved in AC-G115.
There is no need to adjoin \(\sqrt{269}\) and \(\zeta_5\)
separately to the minimal second-branch field.

These fields have Galois group \(S_3\times C_4\) over
\(\mathbb Q\), and degree twenty-four, since the unique
quadratic subfield of the \(S_3\) field is \(K\), whereas
that of \(F_d\) is \(\mathbb Q(\sqrt5)\).
A reflection fixes the degree-twelve field
\(\mathbb Q(\sqrt[3]{269},F_d)\).
Its class group having zero \(269\)-part would suffice,
by the exact reflection norm argument of AC-G103/G126,
to kill the nontrivial cubic target class sectors.
No such new certificate is asserted here.
For a fresh biquadratic host \(K\mathbb Q(\sqrt{-r})
\mathbb Q(\sqrt s)\), the four Shapiro sectors replace
\(d\) by \(d,-rd,sd,-rsd\).
Thus the old degree-six quadratic-character certificates
cannot be substituted into this calculation.

\paragraph{Retained sources and remaining interfaces.}
The integral proof uses precisely the previously audited Katz
modular-curve and q-expansion framework, his full Eisenstein
measure theorem, the complete applied Dasgupta--Kakde constant
formula, Deligne--Serre Lemma 6.11 and Stacks 07S7/023S/023T,
with the changes in \(N\), the tame character and all local
conductors checked above.  The Galois and crystalline comparison
theorems remain published inputs.  The character-line descent
is a genuine finite etale descent and its boundary condition is
vanishing on the full cusp divisor.
An admissible unitary host, native support, exact class-sector
upper bounds, the model/counting comparison and any conclusion
about H04.4 require separate proofs.  No previously fixed
coefficient ring or residual tuple is silently replaced.
